% problem_id: S014-lurie-ultracategories-theorem-2-3-1
% source_id: S014 (Jacob Lurie, Ultracategories)
% source_file: lurie-conceptual.pdf
% statement_locator: printed p. 25
% proof_locator: printed p. [26, 27]
% representation: PDF; transcribed from rendered page images (never from the text layer)
% status: complete (statement + author proof verbatim + reason + locators)
%
%% problem_id: L2
%% theorem_number: Theorem 2.3.1 (Makkai)
%% source: J. Lurie, "Ultracategories"
%% source_locator: /disks/sata1/yupeng/human-proof-corpus/source-cache/registry-sources/lurie/lurie-conceptual.pdf
%%   statement: printed page 25 = PDF page 25
%%   proof:     printed pages 26-27 = PDF pages 26-27
%%
%% Reading method: the PDF was rendered on the host with
%%   pdftoppm -f <p> -l <p> -r 150 -png lurie-conceptual.pdf /tmp/L2-<p>
%% and additionally re-rendered at 300 dpi and cropped (-x/-y/-W/-H) over the lemma
%% statement, the diagram, the tail of the lemma proof, the statement of Theorem 2.3.1
%% and the last lines of its proof. Every line transcribed below was READ FROM THE PAGE
%% IMAGES. The pdftotext text layer of this file is corrupted (it renders the element
%% symbol as "2", alpha as a dagger, and drops the "of Construction 2.2.1" spacing), so
%% it was used ONLY to locate the theorem number and was never treated as source text.
%%
%% Transcription policy: verbatim. Nothing rewritten, simplified or completed; nothing
%% silently corrected. The author's line breaks are not reproduced, but displayed
%% formulas and commutative diagrams are reproduced exactly as displayed. The end-of-proof
%% box is the open square the paper prints.
%%
%% Scope note: the proof of Theorem 2.3.1 proper is the paragraph beginning
%% "Proof of Theorem 2.3.1." on page 26 and ending on page 27. The two numbered items
%% that this proof cites but that the paper states separately --- Theorem 2.2.2 and
%% Corollary 2.1.4 (elsewhere in the paper, not transcribed) and Lemma 2.3.6 (on page 26,
%% immediately before, introduced by the author as "the following observation") --- are
%% not part of the proof of Theorem 2.3.1. Lemma 2.3.6 and its proof ARE transcribed
%% below, inside an explicitly delimited auxiliary block, because the proof of
%% Theorem 2.3.1 delegates its decisive step to it; the block is marked as auxiliary so
%% that it is not mistaken for part of the proof of Theorem 2.3.1.
%% Material on these pages that belongs to neither the statement nor the proof
%% (continuation of Remark 2.3.4, Corollary 2.3.5, the section heading of §2.4) is not
%% transcribed.

\documentclass[11pt]{article}
\usepackage{amsmath}
\usepackage{amssymb}
\usepackage{amsthm}

\theoremstyle{plain}
\newtheorem{theorem}{Theorem}
\newtheorem{lemma}[theorem]{Lemma}
\theoremstyle{remark}
\newtheorem{remark}[theorem]{Remark}

\newcommand{\Mod}{\operatorname{Mod}}
\newcommand{\Set}{\operatorname{Set}}
\newcommand{\Shv}{\operatorname{Shv}}
\newcommand{\Fun}{\operatorname{Fun}}
\newcommand{\ev}{\operatorname{ev}}
\newcommand{\op}{\mathrm{op}}
\newcommand{\Ult}{\mathrm{Ult}}
\newcommand{\LUlt}{\mathrm{LUlt}}

\begin{document}

%% =====================================================================
%% --- page 25 ---
%% =====================================================================

\noindent\textbf{2.3. Application: Strong Conceptual Completeness.} We now turn to the
original version of Makkai's strong conceptual completeness theorem.

\begin{theorem}[Makkai]\label{thm:2.3.1}
Let $\mathcal{C}$ be a small pretopos. Then the evaluation map
$\ev : \mathcal{C} \to \Fun^{\Ult}(\Mod(\mathcal{C}), \Set)$ of
Construction 2.2.1 is an equivalence of categories.
\end{theorem}

\begin{remark}\label{rem:2.3.2}
Theorem 2.3.1 is actually slightly stronger than the version proved by Makkai in [9], since
our category of ultrafunctors $\Fun^{\Ult}(\Mod(\mathcal{C}), \Set)$ is \emph{a priori}
larger than the one introduced by Makkai; see Warning 0.0.9.
\end{remark}

%% =====================================================================
%% --- page 26 ---
%% (The first lines of page 26 finish Remark 2.3.4 and state Corollary 2.3.5; they are
%%  not part of Theorem 2.3.1 or its proof and are not transcribed here.)
%% =====================================================================

We will deduce Theorem 2.3.1 from Theorem 2.2.2 together with the following observation:

%% ==== BEGIN AUXILIARY ITEM (stated and proved by the author on page 26; NOT part of
%% ==== the proof of Theorem 2.3.1, but the proof of Theorem 2.3.1 cites it for its
%% ==== quasi-compactness half. Transcribed verbatim because the proof delegates to it.)

\begin{lemma}\label{lem:2.3.6}
Let $\mathcal{C}$ be a small pretopos and let $F : \Mod(\mathcal{C}) \to \Set$ be a functor
equipped with a left ultrastructure $\{\sigma_\mu\}$. Suppose that, for every collection of
models $\{M_s\}_{s \in S}$ of $\mathcal{C}$ and every ultrafilter $\mu$ on $S$, the map
$\sigma_\mu : F(\int_S M_s d\mu) \to \int_S F(M_s)d\mu$ is surjective. Then $F$ is a
quasi-compact object of the topos $\Fun^{\LUlt}(\Mod(\mathcal{C}), \Set)$.
\end{lemma}

\begin{proof}
Let $\{F_s\}_{s \in S}$ be the collection of all quasi-compact subobjects of $F$ (in the
Grothendieck topos $\Fun^{\LUlt}(\Mod(\mathcal{C}), \Set)$). We regard $S$ as partially
ordered by inclusion (so that $s \leq s'$ if and only if $F_s \subseteq F_{s'}$). Since the
collection of quasi-compact subobjects of $F$ is closed under finite unions, the partial
ordering on $S$ is directed. Applying Proposition 1.1.10, we can choose an ultrafilter
$\mu$ on $S$ such that, for every $t \in S$, we have $\mu(\{s \in S : s \geq t\}) = 1$.

Assume that $F$ is not quasi-compact. Then, for each $s \in S$, we have
$F_s \subsetneq F$. We can therefore choose a model $M_s$ of $\mathcal{C}$ and an element
$x_s \in F(M_s)$ which does not belong to $F_s(M_s)$. Let $x$ denote the image of
$\{x_s\}_{s \in S}$ in the ultraproduct $\int_S F(M_s)d\mu$. Using the surjectivity of the
map $\sigma_\mu : F(\int_S M_s d\mu) \to \int_S F(M_s)d\mu$, we conclude that there is an
element $y \in F(\int_S M_s d\mu)$ satisfying $\sigma_\mu(y) = x$. Then $y$ belongs to
$F_t(\int_S M_s d\mu)$ for some $t \in S$. Using the commutativity of the diagram
\[
\begin{array}{ccc}
F_t(\int_S M_s d\mu) & \xrightarrow{\ \sigma_\mu\ } & \int_S F_t(M_s)d\mu \\[2.2ex]
\downarrow & & \downarrow \\[2.2ex]
F(\int_S M_s d\mu) & \xrightarrow{\ \sigma_\mu\ } & \int_S F(M_s)d\mu,
\end{array}
\]
we see that $x$ can be lifted to an element $\widetilde{x}$ of the ultraproduct
$\int_S F_t(M_s)d\mu$. Choose a subset $S_0 \subseteq S$ satisfying $\mu(S_0) = 1$ and a
tuple $\{\widetilde{x}_s \in F_t(M_s)\}_{s \in S_0}$ representing $\widetilde{x}$. Shrinking
$S_0$ if necessary, we may assume that $\widetilde{x}_s = x_s$ for each $s \in S_0$ and that
$S_0 \subseteq \{s \in S : s \geq t\}$. Then, for any element $s \in S_0$, we conclude that
\[
x_s = \widetilde{x}_s \in F_t(M_s) \subseteq F_s(M_s),
\]
contradicting our choice of $x_s$.
\end{proof}

%% ==== END AUXILIARY ITEM ====

\begin{proof}[Proof of Theorem 2.3.1]
Let $\mathcal{C}$ be a small pretopos; we wish to prove that the evaluation map
$\ev : \mathcal{C} \to \Fun^{\Ult}(\Mod(\mathcal{C}), \Set)$ is an equivalence of categories.
Let $\theta : \Fun^{\LUlt}(\Mod(\mathcal{C}), \Set) \to \Shv(\mathcal{C})$ be the equivalence
of Theorem 2.2.2, so that the composition $\theta \circ \ev : \mathcal{C} \to \Shv(\mathcal{C})$
is the Yoneda embedding. Applying Theorem C.6.5, we see that the evaluation map $\ev$ is a
fully faithful embedding, whose essential image consists of those ultrafunctors
$F : \Mod(\mathcal{C}) \to \Set$ which are quasi-compact and quasi-separated when viewed as
objects of the topos $\Fun^{\LUlt}(\Mod(\mathcal{C}), \Set)$. We will complete the proof by
showing that every ultrafunctor $F$ is quasi-compact and quasi-separated as an object of
$\Fun^{\LUlt}(\Mod(\mathcal{C}), \Set)$. The quasi-compactness of $F$ follows from
Lemma 2.3.6. To prove that $F$ is quasi-separated, we must show that for every pair of

%% =====================================================================
%% --- page 27 ---
%% =====================================================================

quasi-compact objects $F_0, F_1 \in \Fun^{\LUlt}(\Mod(\mathcal{C}), \Set)$ equipped with maps
$F_0 \to F \leftarrow F_1$, the fiber product $F_0 \times_F F_1$ is quasi-compact. It follows
from Theorem 2.2.2 that the topos $\Fun^{\LUlt}(\Mod(\mathcal{C}), \Set)$ is generated (under
small colimits) by objects of the form $\ev_C$, for $C \in \mathcal{C}$. We may therefore
assume without loss of generality that $F_0$ and $F_1$ belong to the essential image of the
evaluation functor, and therefore belong to the subcategory
$\Fun^{\Ult}(\Mod(\mathcal{C}), \Set) \subseteq \Fun^{\LUlt}(\Mod(\mathcal{C}), \Set)$. In this
case, the fiber product $F_0 \times_F F_1$ is also an ultrafunctor (Corollary 2.1.4), and is
therefore quasi-compact by Lemma 2.3.6.
\end{proof}

\end{document}
